% 第 3 章：Parity 全空间 NLP 建模
\section{Parity 全空间 NLP 建模}\label{sec:parity}

\subsection{功能定位与入口}
\compmeta{opf::parity（ParityProblem）}{\srcpath{include/hacdcpf/optimal\_power\_flow/formulation.hpp}}
Parity formulation 是供原生 IPM 与 Ipopt 共用的“可审查”全空间 NLP
（formulation.hpp:11--15）：
\[
\min_x\ f(x)\quad \mathrm{s.t.}\quad g(x)=0,\quad h(x)\le 0,\quad x_{min}\le x\le x_{max}.
\]
装配入口 \srcpath{parity::build\_problem}（src/optimal\_power\_flow/parity\_formulation.cpp:108）。
它是\textbf{唯一同时优化 AC/DC/VSC/DCDC/储能/可切负荷/可再生/柔性负荷/能量路由器的路径}；
混合算例由顶层自动强制启用。

\subsection{变量布局}
\srcpath{VarIndex}（formulation.hpp:66--87；装配 parity\_formulation.cpp:338--406）：
\begin{align*}
x=[\,&\theta;\ V_m;\ P_g;\ Q_g;\ V_{dc};\ P_{ac};\ Q_{ac};\ P_{dc};\ \Delta P_d;\ \Delta Q_d;\
P_{ren};\ Q_{ren};\\
&P_{stor};\ Q_{stor};\ P_{storDC};\ P_{dcdc};\ P_{flex};\ ER_P;\ ER_Q].
\end{align*}
纯 AC 与混合统一；空家族长度 0 但偏移有效。
等式行序（\srcpath{ConstraintIndex}，formulation.hpp:103--122；计数 :408--467）：
$[gP_{ac};\ gQ_{ac};\ gP_{dc};\ g_{vsc};\ g_{dcdc}(=0);\ g_{er};\ g_{dc\_ref}]$
（$n_{dcdc\_bal}$ 仍分配但评估器留空，实现怪癖见 11c 节技术笔记）。
非线性不等式行序：from 端 $|S|^2$、to 端 $|S|^2$、VSC 容量圆、DC 支路 $P^2$、
VSC $I_{ac}^2$、$m$ 上限、$m$ 下限、DC/DC 占空比（每受限台 2 行线性行）。

\subsection{目标函数}
\srcpath{objective}（parity\_formulation.cpp:1138--1240；梯度/Hessian :1242--1329）：
\begin{itemize}
  \item \textbf{Economic}：
        \begin{align*}
        f={}&\sum_g(c_2P_g^2+c_1P_g+c_0)+10^{-3}\sum_c(Q_{ac}-q_{set})^2
        +\mathrm{VOLL}\sum(\Delta P_d+\Delta Q_d)\\
        &{}+\sum \text{固定DER线性成本}
        +\sum_{ren}\big[c_1P+c_{curt}(P_{rated}-P)\big];
        \end{align*}
        无功设定点正则化常数 \texttt{kReactiveSetpointRegularization}$=10^{-3}$（:28，
        注释说明动机：消除平坦方向）；VOLL 自动值
        $\max(100,\ \text{负荷}\ \fld{cost\_mw},\ 1.1\times\max\text{边际成本})$（:64--83）；
  \item \textbf{VoltageDeviation}：$w_v\sum_i(V_i-V^*)^2$；
  \item \textbf{ActiveLoss}：$w_l\,(\sum P_g+P_{ren}+P_{stor}+P_{storDC}-P_{flex}+\Delta P_d)\,S_b$
        （物理损耗 = 发电 $-$ 服务负荷，:1155--1171）；
  \item \textbf{VoltageDeviationAndLoss}：组合目标（:1142--1147，
        RPO Combined 内层使用，reactive\_power\_opt.cpp:344--345）；
  \item \textbf{固定 DER 线性成本是常数项}（不影响极小解，仅影响报告值）：
        AC/DC static gen、不可削减可再生、不可控 PV 的 $c_1\max(0,P\cdot\fld{scaling})$，
        DC PV 阵列用 \fld{p\_set\_mw}（:1200--1216）；
  \item \textbf{可再生削减成本}：仅当 \fld{cost\_curtail\_mwh}$>0$ 且 $p_{rated}-p>0$
        时收取，梯度 $(c_1-\max(0,c_{curt}))\,S_b$（:1231--1233, 1322--1323）；
        注意削减量对 $p_{rated}$ 计量而上界是 $\min(\text{额定},\text{可用出力})$——
        可用$<$额定时差额恒计入成本；PV 系统只有 $c_1$ 无削减价（:1234--1237）；
  \item 非经济目标加 $10^{-6}$ 经济 tie-break（:1177--1183）。
\end{itemize}

\subsection{等式约束}
装配 :1335--1572。

\paragraph{固定注入台账（\fld{p\_fixed\_inj}/\fld{q\_fixed\_inj}，:502--564）}
AC 侧固定注入组成：StaticGenerator（$\times$scaling）、不可削减 RenewableGen、
不可控 PVSystem（经 \srcpath{compute\_pv\_power\_mw} 功率曲线，:533）、
VPP（pcc 出力）、并网 Microgrid（仅 \fld{p\_exchange}）、非 InTransit 的
MobileStorage。DC 侧平衡另有组成：\fld{dc\_loads} 存在时逐负荷
\srcpath{effective\_load\_p\_mw}、dc\_static\_generators、dc\_pv\_arrays 注入
（:1453--1488）。

AC 平衡（ZIP 负荷，\fld{scale\_p}/\fld{scale\_q}
$=1/\max(|P_d|,1)$ 预缩放，:1431--1440）：
\begin{align*}
gP_i={}&P_i^{net}-P_{g,i}-P_{fixed,i}-P_{ren,i}-P_{stor,i}\\
&{}+P_{d,i}(V)-\Delta P_{d,i}+P_{flex,i}-P_{ac,i}=0
\end{align*}
$Q$ 行同构。DC 平衡：
\[
gP_k=p_{d,k}+V_k\sum_m G_{km}V_m-P_{conv\_dc,k}=0;
\]
DC/DC 折叠进 DC 行：输入侧 $+p_{in}$，输出侧 $-p_{out}$，
$p_{out}=p_{in}\eta$（正向）或 $p_{in}/\eta$（反向），加 $I^2R$ 损耗
$r_{eq}p_{in}^2/v_{dc,in}^2$（:1497--1514）。
VSC 耦合：$P_{ac}+P_{dc}+a+bI_{ac}+cI_{ac}^2=0$（:1516--1530，$\varepsilon=\fld{eps\_iac}=10^{-6}$）。
能量路由器：端口功率为决策变量，每路由器一条\textbf{无损}守恒行 $\sum P_{port}=0$
（:1548--1566；注释明确 \fld{loss\_percent} 细化 deferred，:1563--1564）。
DC\_V 参考电压行 $V_{dc,k}-v_{set}=0$（:1567--1571）；
无导电支路的 DC 母线通过 $\pm10^{-6}$ 界锚定（:613--627）。

\subsection{不等式约束}
$h(x)\le0$ 约定，:1824--1925：
\begin{align*}
\text{AC 支路两端：}\quad &P_f^2+Q_f^2-S_{max}^2\le0\quad(\pi\text{ 模型，含 tap 与 \fld{shift\_deg}，:1836--1861})\\
\text{VSC 容量圆：}\quad &P_{ac}^2+Q_{ac}^2-S_{max}^2\le0\quad(:1862--1867)\\
\text{DC 支路：}\quad &P_{km}^2-P_{max}^2\le0,\quad P=gV_k(V_k-V_m)\quad(:1869--1879)\\
\text{VSC 电流限：}\quad &P_{ac}^2+Q_{ac}^2-i_{ac,max}^2V_m^2\le0\quad(:1881--1890)\\
\text{调制度（线性行）：}\quad &m=\frac{V_mV_{n,ac}}{K_mV_{dc}V_{n,dc}}\ \text{上下限}\quad(:1891--1908)\\
\text{DC/DC 占空比（线性行）：}\quad &\text{见下方全拓扑公式表}\ (:1797--1818,\ 1909--1924)
\end{align*}
注意口径：AC 支路限为 $|S|$（视在功率）双端圆约束；DC 支路限为 $|P|$。

\paragraph{换流器容量圆半径 \srcpath{converter\_smax\_pu}（:85--101）}
优先 \fld{conv.p\_rated\_mw}，无效时回退
$\max(|p_{max}|,|p_{min}|,|q_{max}|,|q_{min}|,10^{-6})$；
\fld{enforce\_converter\_capacity}=false 时返回 $10^{12}$ 使圆恒不激活。

\paragraph{准入条件（未声明限值的设备不占行，:435--464）}
$I_{ac}$ 行需 \fld{isfinite(i\_ac\_max\_pu)}$>0$；$m$ 行需
\fld{k\_m\_modulation}$>0\ \land\ v_{n,ac}>0\ \land\ v_{n,dc}>0\ \land\ m_{max}>0$
（$m_{min}$ 行同理）；占空比行需 topology$\ne$Generic 且 $0\le d_{min}<d_{max}\le1$；
端口额定 kV 取 \fld{vn\_in/out\_kv}，缺省回退母线 base\_kv 再回退 1.0（:1791--1795）。

\paragraph{DC/DC 占空比行全拓扑公式（\srcpath{dcdc\_duty\_rows}，:1797--1818）}
每受限台 2 行 $h=a_{in}V_{in}+a_{out}V_{out}\le0$（上限行 $D\le d_{max}$、
下限行 $D\ge d_{min}$），$n$=\fld{n\_ratio}（默认 1）：
\begin{center}
\footnotesize
\begin{tabular}{@{}llll@{}}
\toprule
拓扑 & 电压律 & $D\le d_{max}$ 行 $(a_{in},a_{out})$ & $D\ge d_{min}$ 行\\
\midrule
Buck & $V_{out}=DV_{in}$ & $(-d_{max}Vn_{in},\ Vn_{out})$ & $(d_{min}Vn_{in},\ -Vn_{out})$\\
Boost & $V_{out}=V_{in}/(1-D)$ & $(-Vn_{in},\ (1-d_{max})Vn_{out})$ & $(Vn_{in},\ -(1-d_{min})Vn_{out})$\\
BuckBoost & $V_{out}=\tfrac{D}{1-D}V_{in}$ & $(-d_{max}Vn_{in},\ (1-d_{max})Vn_{out})$ & $(d_{min}Vn_{in},\ -(1-d_{min})Vn_{out})$\\
Isolated & $V_{out}=nDV_{in}$ & $(-d_{max}nVn_{in},\ Vn_{out})$ & $(d_{min}nVn_{in},\ -Vn_{out})$\\
\bottomrule
\end{tabular}
\end{center}

\subsection{变量界与初值}
界（:569--754）：$\theta\in[-\pi,\pi]$；$V_m$ 母线限；$P_g/Q_g$ 机组限
（退役机组钉死 $\pm10^{-8}$）；储能钉在调度值 $\pm10^{-4}/S_b$
（无跨时段 SOC 模型的诚实说明，:674--717）；\textbf{例外}：
\fld{cap\_charging\_strategy}==``static'' 时 $P$ 钉在 $-$\fld{cap\_static\_charging\_mw}、
$Q$ 钉 0（$\pm10^{-4}/S_b$ 容差带），即承载力静态 ESS 完全退出优化；
其余储能钉调度值 \fld{p\_mw}$\pm$eps 且同时受 $[p_{min},p_{max}]$ 夹（:679--717）；
可再生上限 $=\min(\text{额定},\text{可用出力})$
（:651--672）；柔性负荷 $[p-\fld{flex\_down},\ p+\fld{flex\_up}]$（:730--734）；
DC/DC 界 $p_{max}=\fld{pmax\_mw}>10^{-9}\,?\,\fld{pmax\_mw}:\fld{sn\_mva}$，
$p_{min}\ge p_{max}$ 时取 $-p_{max}$（:719--727）。

初值链：母线设定点 $\to$ 零出力机组按 $P^{max}$ 比例分摊缺额（:940--966）
$\to$ DC 潮流暖启动（稀疏 $B'$ 求解，:757--857）$\to$ $Q_g$ 暖启动（:861--917）
$\to$ 切负荷吸收残差（:1038--1095）$\to$ $J_gJ_g^{\mathsf T}$ 最小二乘可行性修正
（阻尼 0.8，:1097--1126）$\to$ 1\% 内点钳制。
DC/DC 暖启动 \fld{p\_ref} 是输出参考，经 $\eta$（clamp $[0.01,1]$）换算成输入参考
（:1004--1012）；ER 端口注入在 AC 行缩放之后补乘 \fld{scale\_p}/\fld{scale\_q}
（:1545--1556）。

\subsection{解析导数}
解析等式 Jacobian（:1574--1780）、不等式 Jacobian（:1927--2071）、
稀疏 Lagrangian Hessian $\nabla^2L=\nabla^2f+\sum\lambda\nabla^2g+\sum\nu\nabla^2h$
（:2087--2560，含 ZIP 二阶项、VSC 损耗全二阶、DC/DC $I^2R$ 二阶、支路 $|S|^2$ 全二阶）。

\paragraph{ZIP 系数与二阶导}
系数：有组件负荷时用 SolverData 逐母线加权 \fld{bus\_zip\_*}，否则全局
\fld{zip\_pw}/\fld{zip\_qw}（默认恒功率）（:133--160）；
Hessian 中 ZIP 二阶项 $2w_2p_d$（P 行）、$2w_2q_d$（Q 行）加在 $V_m$ 对角
（:2235--2239），且 AC 平衡行曲率乘 $\lambda\cdot\fld{scale\_p}/\fld{scale\_q}$
（:2187--2188）。

\paragraph{Hessian 稀疏结构与关键二阶项}
下三角 triplet 双边写入、去重求和装配；$|\lambda|$ 或 $|\nu|<10^{-14}$ 的行整块跳过
（:2189 等）；调制/占空比行线性、零曲率，仅用游标推进保持对偶对齐（:2548--2553）。
DC/DC $I^2R$ 项（$g_{out}$ 中 $+rp_{in}^2/v^2$）：
$\partial^2/\partial p^2=2r/v^2$、$\partial^2/\partial p\partial v=-4rp/v^3$、
$\partial^2/\partial v^2=6rp^2/v^4$（:2319--2351，$v$ 取 $\max(v_{in},0.1)$）；
$I_{ac}$ 限行 $\partial^2h/\partial V_m^2=-2i_{max}^2$（:2546）。

正确性校核：FD 诊断 \srcpath{check\_ac\_opf\_jacobian\_fd}（死代码）与测试中的
FD 断言（见第~\ref{sec:validation}~章）。

\subsection{参数表（ParityOptions）}
定义 formulation.hpp:35--49。
\begin{paramtable}
\fld{load\_shedding} & -- & bool & true/false & true & 启用可切负荷变量 $\Delta P_d,\Delta Q_d$\\
\fld{voll} & VOLL & \$/MWh & $\ge100$ & 0（自动） & 切负荷价值；0 = $\max(100,1.1\times\max$边际成本$)$\\
\fld{eps\_iac} & $\varepsilon$ & pu & $10^{-8}$--$10^{-4}$ & $10^{-6}$ & VSC 电流数值正则化\\
\fld{objective} & -- & 枚举 & 四档 & Economic & Economic / VoltageDeviation / ActiveLoss / VoltageDeviationAndLoss\\
\fld{voltage\_target\_pu} & $V^*$ & pu & 0.95--1.05 & 1.0 & 电压偏差目标\\
\fld{voltage\_deviation\_weight} & $w_v$ & -- & $\ge0$ & 1.0 & 电压偏差目标权重\\
\fld{active\_loss\_weight} & $w_l$ & -- & $\ge0$ & 1.0 & 网损目标权重\\
\fld{enforce\_branch\_limits} & -- & bool & -- & true & AC 支路 $|S|$ 双端约束\\
\fld{enforce\_converter\_capacity} & -- & bool & -- & true & VSC 容量圆\\
\fld{enforce\_converter\_current\_limits} & -- & bool & -- & true & VSC $I_{ac}$ 限\\
\fld{enforce\_converter\_modulation\_limits} & -- & bool & -- & true & 调制度 $m$ 与 DC/DC 占空比行\\
\end{paramtable}

\subsection{验证与校核}
\begin{itemize}
  \item parity 成本不劣于经济调度 $\times1.001+1$（test\_opf\_solver\_backends.cpp:537--553）；
  \item 夜间纯 external-grid 算例：储能钉调度值 WithinAbs $2\times10^{-4}$、
        DC/DC $I^2R$ 公式复核、AC 供需平衡 $\pm10^{-3}$
        （test\_nighttime\_opf.cpp:46--96）；
  \item DC slack KCL Jacobian FD $<10^{-6}$（test\_opf\_solver\_backends.cpp:267--301）；
  \item 结果中的 \fld{converter\_model\_scope} 按实际启用约束族拼装标签
        （ac\_opf.cpp:2254--2272）。
\end{itemize}
